\documentclass[12pt]{article}
\everymath{\displaystyle}
\usepackage{unicode-math}
\usepackage{pgfplots}
\pgfplotsset{compat=1.18}
\usepackage{amsmath}


\usepackage[margin=1in]{geometry}
\usepackage{tcolorbox}
\usepackage{graphicx}
\usepackage{tikz}
\usepackage{amssymb}
\usepackage{enumitem}
\usepackage{titlesec}
\usepackage{fancyhdr}
\usepackage{pifont}
\usepackage{xcolor}
\newcounter{questionnum}
\setcounter{questionnum}{0}
\usepackage{eso-pic} % for page background

% --------- 主题颜色设置 ---------
\definecolor{novelblue}{RGB}{70,60,150}
\definecolor{headerbg}{RGB}{240,240,255}

\pagestyle{fancy}
\fancyhf{}
\renewcommand{\headrulewidth}{0.4pt}
\renewcommand{\headrule}{\hbox to\headwidth{\color{novelblue}\leaders\hrule height 0.4pt\hfill}}

% --------- 页眉背景色条 ---------
\newcommand{\headbackground}{
  \AddToShipoutPictureBG*{
    \begin{tikzpicture}[remember picture, overlay]
      \fill[blue!5!purple!10] (current page.north west) rectangle ([yshift=-60pt]current page.north east);
    \end{tikzpicture}
  }
}

% --------- 页眉设置 ---------
\pagestyle{fancy}
\fancyhf{}
\renewcommand{\headrulewidth}{0pt}
\fancyhead[L]{\color{novelblue}\textbf{\textbf{Novel Prep} }}
\fancyhead[C]{\color{novelblue}\textbf {SAT Preparation}}
\fancyhead[R]{\color{novelblue}\textbf{Jack Zeng}}


\newtcolorbox{SATCard}[1][]{%
  enhanced,
  colback=white,
  colframe=novelblue,
  coltitle=white,
  boxrule=0.6pt,
  arc=4pt,
  boxsep=5pt,
  left=8pt, right=8pt, top=10pt, bottom=10pt,
  sharp corners=south,
  drop shadow south east,
  #1
}

% --------- 顶部 Mark for Review 栏 ---------
\newcommand{\markforreview}{
  \stepcounter{questionnum} % 自动递增题号
  \begin{tikzpicture}[scale=1.1]
    % 黑色题号
    \fill[black] (0,0) rectangle (0.8,0.6);
    \node[white, font=\bfseries\large] at (0.4,0.3) {\thequestionnum};

    % 灰底主栏
    \fill[gray!10] (0.8,0) rectangle (14.5,0.6);

    % Mark for Review
    \node[anchor=west, font=\small] at (1.2,0.3) {\ding{72} \hspace{2pt}Mark for Review};

    % ABC 按钮区域
    \draw[thick, rounded corners=3pt] (13.5,0.12) rectangle (14.5,0.48);
    \node[font=\bfseries\normalsize] at (14.0,0.3) {ABC};
  \end{tikzpicture}


  \newcommand{\choiceboxcircled}[2]{%
  \begin{tcolorbox}[
    colback=white,
    colframe=black,
    boxrule=0.5pt,
    rounded corners=6pt,
    left=6pt, right=6pt, top=6pt, bottom=6pt,
    width=\linewidth
  ]
    \textbf{\large\textcircled{\textbf{#1}}} \ #2
  \end{tcolorbox}
}
}

% --------- 选项框样式 ---------
\newcommand{\choicebox}[2]{%
  \begin{tcolorbox}[
    colback=white, 
    colframe=novelblue,
    boxrule=0.5pt, 
    rounded corners=4pt, 
    left=6pt, right=6pt, top=6pt, bottom=6pt,
    width=\linewidth
  ]
  \textbf{#1.} \ #2
  \end{tcolorbox}
}


\begin{document}
\section*{Unit 2 Advanced Math Practice Test Module II}

\noindent\markforreview
\vspace{1em}

For what value of \(x\) is the given expression equivalent to \((41k)^{25x}\), where \(k > 1\)?

\[
\sqrt[8]{41k} (\sqrt[9]{41k} )^{12}
\]

\vspace{0.75em}

\vspace{2em}

\noindent\markforreview
\vspace{1em}

The expression \(4^{18x}\) is equivalent to \(k^{3x}\), where \(k\) is a constant.  
What is the value of \(k\)?

\vspace{0.75em}

\vspace{2em}

\noindent\markforreview
\vspace{1em}

The expression \(\frac{(6x + 8)(3x + 6)}{12x + 24}\) is equivalent to \(\frac{ax + b}{2}\), where \(a\) and \(b\) are constants and \(x > 0\).  
What is the value of \(a + b\)?

\vspace{0.75em}

\vspace{2em}

\noindent\markforreview
\vspace{1em}

Which expression is equivalent to \(\left(3\sqrt{a} - \sqrt{b}\right)^{\frac{2}{7}}\), where \(a > b\) and \(b > 0\)?

\choicebox{A}{\(\sqrt[7]{9a-b} \)}  
\choicebox{B}{\(\sqrt[7]{9a- 6ab + b} \)}  
\choicebox{C}{\(\sqrt[7]{9a- 3\sqrt{ab} + b} \)}  
\choicebox{D}{\(\sqrt[7]{9a- 6\sqrt{ab} + b} \)}  

\vspace{2em}
\newpage
\noindent\markforreview
\vspace{1em}

The expression \(12x^2 - mx - 20\), where \(m\) is a constant, is equivalent to \((kx - 5)(nx + 4)\), where \(k\) and \(n\) are integer constants.  
Which of the following must be an integer?

\choicebox{A}{\(\frac{12}{m}\)}  
\choicebox{B}{\(\frac{12}{n}\)}  
\choicebox{C}{\(\frac{m}{k}\)}  
\choicebox{D}{\(\frac{m}{n}\)}  

\vspace{2em}

\noindent\markforreview
\vspace{1em}

In the given equation, \(m\) and \(n\) are constants and \(n < 0\).  
If the equation is true for all values of \(x\), what is the value of \(m + n\)?

\[
4x^2 + mx + 9 = (2x + n)^2
\]

\choicebox{A}{\(-15\)}  
\choicebox{B}{\(-9\)}  
\choicebox{C}{\(-3\)}  
\choicebox{D}{\(12\)}  

\vspace{2em}
\newpage
\noindent\markforreview
\vspace{1em}

The graph of \(y = f(x)\) has a \(y\)-intercept at \((0, 12)\) and \(x\)-intercepts at \((-3, 0)\) and \((2, 0)\).  
Which of the following equations could define \(f\)?

\choicebox{A}{\(f(x) = (x + 3)^2(x - 2)\)}  
\choicebox{B}{\(f(x) = (x + 3)(x - 2)^2\)}  
\choicebox{C}{\(f(x) = (x - 3)^2(x + 2)\)}  
\choicebox{D}{\(f(x) = (x - 3)(x + 2)^2\)}  

\vspace{2em}

\noindent\markforreview
\vspace{1em}

Function \(f\) is defined by \(f(x) = a\sqrt{x} + b\), where \(a\) and \(b\) are constants.  
In the \(xy\)-plane, the graph of \(y = f(x)\) has an \(x\)-intercept at \((16, 0)\), and the graph of \(y = f(x) - 10\) has an \(x\)-intercept at \((9, 0)\).  
What is the value of \(b\)?

\vspace{2em}

\noindent\markforreview
\vspace{1em}

For the function \(f(x) = b\sqrt{a - x} + 7\sqrt{a - x}\),  
\(a\) and \(b\) are constants and \(f(16) = 0\).  
If \(f(9) > 0\), which of the following could be the value of \(a + b\)?

\choicebox{A}{\(-3\)}  
\choicebox{B}{\(7\)}  
\choicebox{C}{\(9\)}  
\choicebox{D}{\(15\)}  

\vspace{2em}
\newpage
\noindent\markforreview
\vspace{1em}

What is the sum of the solutions to the given equation?

\[
x^2 + 7x + 19 = 28 - 5x
\]

\choicebox{A}{\(-12\)}  
\choicebox{B}{\(-9\)}  
\choicebox{C}{\(-7\)}  
\choicebox{D}{\(12\)}  


\vspace{2em}

\noindent\markforreview
\vspace{1em}

What are the solutions to the given equation?

\[
4x^2 - 9x - 1 = 0
\]

\choicebox{A}{\(x = \frac{-9 \pm \sqrt{97}}{8}\)}  
\choicebox{B}{\(x = \frac{-9 \pm \sqrt{65}}{8}\)}  
\choicebox{C}{\(x = \frac{9 \pm \sqrt{97}}{8}\)}  
\choicebox{D}{\(x = \frac{9 \pm \sqrt{65}}{8}\)}  

\vspace{2em}
\newpage
\noindent\markforreview
\vspace{1em}

One solution to the given equation can be written as \(\frac{\sqrt{k} - 1}{2}\), where \(k\) is a constant.  
What is the value of \(k\)?

\[
x^2 + x - 5 = 0
\]

\vspace{2em}

\noindent\markforreview
\vspace{1em}

In the given equation, \(k\) is a positive constant.  
What is the product of the solutions to the given equation?

\[
x^2 - 8kx - (k - 4) = 0
\]

